\begin{helper}
Fix an actual paired mixed surface
\(\widehat S_j(\omega),\widehat C_\omega,\lambda,M_j(\omega)\) from
\(\noderef{SamplingSplitOutsideBlockerPairedMixedMeasureSurface}\), and fix
the marginal-preserving coupling data
\(\mathcal C\) from
\(\noderef{SamplingSplitOutsideBlockerMarginalPreservingCouplingData}\) on
the same atom family \(C_\omega,C'_\omega,w(\omega)\).  Assume that the
surface's paired measure and paired cells are exactly the paired measure and
paired cells of \(\mathcal C\).  Assume also that \(\phi\) is injective, that
the actual list \(z_0,\ldots,z_{n-1}\) covers every source outside blocker
adjacent to \(X\), that the source law \(\nu\) satisfies the Bernoulli
activation atom formula and priority lower-rectangle law with a parameter
\(p\in[0,1]\), and that the source endpoint event
\[
  E_X=\{(A,\pi): I_G(A,\pi)\cap X\ne\emptyset\}
\]
is measurable.  Then for every atom \(\omega\),
\[
  \operatorname{ofReal}(w(\omega)M_0(\omega))
  =
  \nu(E_X\cap C_\omega).
\]
\end{helper}
\begin{proof}
Fix \(\omega\).  The measured-stage field of
\(\noderef{SamplingSplitOutsideBlockerPairedMixedMeasureSurface}\), applied
with \(j=0\), gives
\[
  \operatorname{ofReal}(w(\omega)M_0(\omega))
  =
  \lambda(\widehat S_0(\omega)\cap\widehat C_\omega).
\]
We first compare \(\widehat S_0(\omega)\), inside the paired cell, with the
first-coordinate cylinder over the displayed source local/blocker endpoint
event.  Let \(D_X\) be the event that there is \(x\in X\) such that \(x\) is
active, \(x\) beats every active neighbor lying in \(X\cup\{r\}\), and every
listed blocker \(z_i\) with \(x\in B_i\) fails the opposite strict test
\[
  \neg(z_i\in A,\ \pi(x)<\pi(z_i),\ 0\le \pi(z_i)\le1).
\]
This is only the displayed endpoint event; it is not yet identified with the
true endpoint survival event.

Let \(\zeta=((A,\pi),(A',\pi'))\in\widehat C_\omega\).  Expanding the
mixed-stage definition in
\(\noderef{SamplingSplitOutsideBlockerPairedMixedMeasureSurface}\) at
\(j=0\), membership in \(\widehat S_0(\omega)\) means that there is \(x\in X\)
with \(x\in A\) and \(\phi(x)\in A'\), the source and target embedded-neighbor
priority tests hold for vertices in \(X\cup\{r\}\), and every retained outside
blocker \(z_i\) with \(x\in B_i\) satisfies the displayed negated source
blocker test.  There are no private target blocker tests at \(j=0\).

The additional target embedded requirements are redundant on
\(\widehat C_\omega\).  The surface records \(X\cup\{r\}\subseteq S\), where
\(S\) is the source common coordinate set, and on \(\widehat C_\omega\) it
synchronizes common activation traces and priorities.  Since \(\phi\) is
injective, \(x\in A\) and the source embedded comparisons are equivalent on
the paired cell to \(\phi(x)\in A'\) and the corresponding target embedded
comparisons.  Therefore, on \(\widehat C_\omega\),
\[
  \zeta\in\widehat S_0(\omega)
  \quad\Longleftrightarrow\quad
  \zeta_1\in D_X .
\]
Consequently
\[
  \widehat S_0(\omega)\cap\widehat C_\omega
  =
  \{\zeta:\zeta_1\in D_X\}\cap\widehat C_\omega .
\]

The event \(D_X\) is measurable by
\(\noderef{SamplingSplitOutsideBlockerSourceDisplayedEndpointMeasurable}\):
it is a finite union, over \(x\in X\), of finite intersections of activation
cylinder conditions, strict priority inequalities between named source
coordinates, and closed unit-interval conditions on named source coordinates.
Thus the displayed event is measurable before any marginal law is used.

The surface and \(\mathcal C\) have the same paired measure and paired cells.
Applying the source restricted marginal law from
\(\noderef{SamplingSplitOutsideBlockerMarginalPreservingCouplingData}\) to
this measurable displayed event \(D_X\) gives
\[
  \lambda(\{\zeta:\zeta_1\in D_X\}\cap\widehat C_\omega)
  =
  \nu(D_X\cap C_\omega).
\]

It remains to replace \(D_X\) by the true endpoint event \(E_X\) in measure
on the same cell.  Apply
\(\noderef{SamplingSplitOutsideBlockerEndpointBoundaryCellReduction}\) to the
source graph, with candidate type \(\{x:V:x\in X\}\), candidate set the full
finite subtype, embedding \(x\mapsto x\), local coordinate set
\(X\cup\{r\}\), blocker \(b_i(x)=z_i\), and blocker predicate
\(H_i(x)\) equal to \(x\in B_i\).  The local-self condition follows from
\(x\in X\).  If \(H_i(x)\), then the candidate-set field of the paired
surface gives \(x z_i\in E(G)\), so the listed blocker is a genuine neighbor.
For neighbor coverage, let \(y\) be adjacent to a candidate \(x\).  If
\(y\in X\cup\{r\}\), it is local.  Otherwise the outside-blocker coverage
hypothesis gives an index \(i\) with \(z_i=y\), and the surface's
candidate-set identity gives \(x\in B_i\).  Thus the hypotheses of the
boundary-reduction node are exactly the source activation atom law, the source
priority rectangle law, and the listed-cover facts assumed here.  It yields
\[
  \nu(E_X\cap C_\omega)=\nu(D_X\cap C_\omega).
\]
Combining this mass equality with the measured-stage identity and the
restricted source marginal identity proves the claimed source endpoint mass
identity for \(\omega\).  Since \(\omega\) was arbitrary, the result holds for
every atom.
\end{proof}
